\section{Sémantická analýza}\label{sec:semanticka-analyza}

Parser zaručuje,~že program odpovídá gramatice. To ještě neznamená,~že dává smysl. Příkaz \texttt{let x = ;} je syntakticky chybný,~protože za rovnítkem chybí výraz. Příkaz \texttt{print y;} je naopak syntakticky správný,~ale není možné jej vykonat,~pokud proměnná \texttt{y} nebyla deklarována.

\emph{Sémantický analyzátor} proto kontroluje podmínky,~které nelze nebo není vhodné vyjadřovat samotnou gramatikou. Patří mezi ně zejména existence deklarací,~viditelnost jmen,~typová správnost výrazů a~platnost jednotlivých operací.

\subsection{Tabulka symbolů}

Překladač si pro každý deklarovaný identifikátor uchovává informace v~\emph{tabulce symbolů}. U~proměnné to může být jméno,~typ,~oblast platnosti a~místo uložení; u~funkce navíc typy parametrů a~návratová hodnota.

Při čtení deklarace překladač přidá nový záznam. Při každém použití identifikátoru hledá odpovídající záznam v~aktuální a~okolních oblastech platnosti. Díky tomu odhalí použití nedeklarované proměnné i~nepovolenou opakovanou deklaraci.

\subsection{Typová kontrola}

Typová kontrola přiřazuje typy listům stromu a~postupuje směrem ke kořeni. Je-li například \texttt{steps} číslo a~\texttt{speed} číslo,~výsledek jejich násobení je opět číslo. Součet s~číselnou proměnnou \texttt{position} proto může být přiřazen do další číselné proměnné.

\begin{figure}[h]
    \centering
    \begin{tikzpicture}[
    expr/.style={task,fill=blue!8,minimum width=28mm,font=\small},
    value/.style={task,fill=orange!12,minimum width=24mm,font=\small}
]
    \node[value] (steps) at (-3.2,0) {\texttt{steps}\\\texttt{number}};
    \node[value] (speed) at (0,0) {\texttt{speed}\\\texttt{number}};
    \node[expr] (mul) at (-1.6,-1.5) {$*$\\výsledek: \texttt{number}};
    \node[value] (position) at (3.2,-1.5) {\texttt{position}\\\texttt{number}};
    \node[expr] (add) at (0.8,-3.1) {$+$\\výsledek: \texttt{number}};
    \node[value,fill=green!10] (target) at (0.8,-4.7)
        {\texttt{newPosition}\\\texttt{number}};

    \draw[diredge] (steps)--(mul);
    \draw[diredge] (speed)--(mul);
    \draw[diredge] (mul)--(add);
    \draw[diredge] (position)--(add);
    \draw[diredge] (add)--node[right,font=\small]{přiřazení}(target);
\end{tikzpicture}

    \caption{Zjednodušená typová kontrola výrazu.}
    \label{fig:typova-kontrola}
\end{figure}

Jiný jazyk může některé kombinace automaticky převádět,~například celé číslo na desetinné,~zatímco další kombinace odmítne. Pravidla typového systému jsou proto součástí definice konkrétního jazyka,~nikoliv univerzální vlastností všech překladačů.

\paragraph{Sémantická kontrola malého jazyka.}

Náš jazyk má jediný číselný typ. Nemusíme tedy porovnávat různé typy,~ale stále musíme zkontrolovat deklarace. Analyzátor v~programu \ref{lst:semantic} prochází příkazy v~pořadí,~ukládá deklarovaná jména do množiny a~rekurzivně kontroluje výrazy.

\lstinputlisting[
    caption={Sémantická kontrola deklarací a~použití proměnných.},
    label={lst:semantic}
]{04-kompilatory/code/mini_compiler/semantic.py}

Inicializační výraz se kontroluje ještě před vložením nového jména do tabulky. Program \texttt{let x = x + 1;} proto nemůže použít právě deklarovanou proměnnou,~dokud pro ni neexistuje předchozí hodnota. Po úspěšné kontrole ví další fáze,~že každé jméno odkazuje na existující proměnnou.

\tipbox{Syntaktický strom může být po sémantické analýze doplněn o~typy,~odkazy do tabulky symbolů a~další atributy. Hovoříme pak o~anotovaném stromu. Další fáze díky tomu nemusí stejné informace zjišťovat znovu.}
